\documentclass[]{article}

\usepackage[
	nomarginpar,
	left=0.5in,
	right=0.5in,
	top=0.25in,
	bottom=0in,
]{geometry}
\usepackage{
	array,
	colortbl,
	enumitem,
	lipsum,
	parskip,
	microtype,
	ulem,
	xcolor,
}
\usepackage[english]{babel}
\usepackage[T1]{fontenc}
\usepackage[hidelinks]{hyperref}
\usepackage[utf8]{inputenc}
\usepackage[compact]{titlesec}

\setlist[itemize]{leftmargin=*, topsep=0pt, nosep,label={\textbullet}}

\newcolumntype{L}{>{\raggedleft}p{0.13\textwidth}}
\newcolumntype{R}{p{0.8\textwidth}}
\definecolor{lightgray}{gray}{0.8}
\newcommand\VRule{\color{lightgray}\vrule width 0.5pt}
\setlength{\arrayrulewidth}{0.5pt}
\arrayrulecolor{lightgray}

\titleformat{\section}
	{\normalfont\large\bfseries\scshape}
	{\thesection}
	{1em}
	{}
	[{\titlerule[0.5pt]}]

\titleformat{\subsection}
  {\normalfont\normalsize\bfseries}
  {\thesubsection}
  {1em}
  {}

\titlespacing*{\subsection} {0pt}{2ex}{0ex}

\def\name{John D. Duncan, III}
\def\phone{(240) 688-7187}
% \def\email{duncanjdiii+jobs@googlemail.com}
\def\email{cv@johnduncan.io}
\def\webpage{https://johnduncan.io}
\def\linkedin{https://linkedin.com/in/johndduncan}
\def\github{https://github.com/johndduncaniii}

\def\Plus{\texttt{+}}
\def\Minus{\texttt{-}}

\setlength\parindent{0em}
\markright{\name}
\thispagestyle{empty}
\frenchspacing

\begin{document}

{\huge \bfseries \name}

\begin{tabular}{@{}l!{\VRule}l@{}}
	Phone & \phone\\
	Location & Washington, DC\\
	Status & US Citizen \\
\end{tabular}
\hfill
\begin{tabular}{@{}l!{\VRule}r@{}}
	Email & \href{mailto:\email}{\tt \email}\\ % \faGoogle
	Homepage & \href{\webpage}{\tt \webpage}\\
	Github & \href{\github}{\tt \github}\\ % \faGithub
	Linkedin  & \href{\linkedin}{\tt \linkedin}\\ % \faLinkedin
\end{tabular}

Hi Will,

I'm John Duncan, a software engineer who helped scale Quorum Analytics—the best-in-class civic tech and public affairs platform—from \$9M to \$64M+ ARR over nearly seven years. As technical lead, I accelerated growth by solving Quorum's most complex engineering challenges, mentoring dozens of engineers, and consistently delivering major projects on time and under budget. I'm an inquisitive engineer who thrives on solving ambiguous problems, balancing pragmatic trade-offs while maintaining high standards. My natural curiosity drives me to deeply understand systems, leading to elegant and maintainable solutions that can be easily documented, maintained, extended, or deleted.

During my tenure at Quorum, I authored 5,000{\Plus} commits, reviewed 1,000{\Plus} pull requests, and conducted 100{\Plus} interviews. Most recently, I launched PAC Hybrid, expanding Quorum’s business into the \$15M{\Plus} PAC market, where I:
\begin{enumerate}
	\item Discovered critical discrepancies in Federal Election Commission (FEC) national party committee C-number ids across 50{\Plus} years of advisory opinion legal records, website data, and official GitHub repository.
	\item Architected end-to-end PAC search analytics platform spanning four datasets, implementing dozens of optimized Django SQL query methods, RESTful APIs, and interactive data visualizations.
	\item Designed a database synchronization infrastructure to instantly sync billions of rows between Quorum and legacy CisionPAC databases.
	\item Reduced database query latency from 60 seconds to 20 milliseconds by optimizing Postgres partition and schema constraints, improving transaction dataset query performance by over 3000x and permanently fixing a 10{\Plus} year old bug.
	\item Coached engineers to optimize a database migration by replacing sequential queries with a single constant-time operation, reducing execution time from two weeks to under five seconds and improving backfill migration performance by over 240,000x.
\end{enumerate}

I wrote, designed, led, and/or facilitated the PAC Hybrid, Quorum PAC, Search, Dashboards, Committee \& Presidential Transcripts, Bulk Actions, International, and Regulations projects. My full \href{https://johnduncan.io/cv.pdf}{CV} with a detailed project history and \href{https://johnduncan.io/portfolio#quorum}{portfolio} of my work are both available on my homepage: \href{https://johnduncan.io/}{https://johnduncan.io/}.

At Quorum, I developed strong expertise in all aspects of full-stack development:\\
Back-end: Python, Django, PostgreSQL, PostGIS, Elasticsearch, Tastypie\\
Front-end: React, React Native, TypeScript, JavaScript, Redux, D3, HTML5, CSS3

I’ve worked as a software engineer while also collaborating cross-functionally with design and product teams, stepping in as needed to drive progress. For example, I collaborated with our former head of design to implement Quorum's design system and perform user interviews to better understand customer needs while also working closely with the PAC product team to define requirements, write technical design documents (TDDs) and product requirements documents (PRDs), perform legal and industry research, and set overall team deliverables. I played a major role in the overall aesthetics, user interface, and user experience of the Quorum platform over my nearly seven years at the company.

In college, I taught myself JavaScript, HTML, and CSS in order to build my old homepage, which was a 90’s-themed net art project that I worked on between 2014--2017: \href{http://cs.gettysburg.edu/~duncjo01/}{http://cs.gettysburg.edu/\~duncjo01/}.

I originally discovered the MTA's front-end engineering role from your post on Hacker News in July 2024, but the timing wasn't right; now, as I am planning on relocating from DC to New York, I'm excited to apply for the position. Beyond my technical qualifications, I have a strong hobbyist interest im urban transit and city planning—I closely follow developments in mass transit systems \& urban zoning policy and regularly interface with the DC government as a board member of the tallest privately owned building in DC.

I graduated from Gettysburg College with a Bachelor of Science in Computer Science (3.60, Honors) and Philosophy (3.73, Honors) and was awarded the class of 2017 Outstanding Computer Science Student award.

I look forward to hearing back from you!

Thanks,

John

\end{document}
